\documentclass[11pt]{amsart}
\usepackage[margin=1in]{geometry}
\usepackage{amsmath,amssymb,amsthm}
\usepackage[hidelinks]{hyperref}

\newtheorem{theorem}{Theorem}
\newtheorem{lemma}{Lemma}

\DeclareMathOperator{\Iso}{Iso}
\DeclareMathOperator{\Aut}{Aut}

\title{Literature-Implied Answer to Question \texttt{quest:kpt} in arXiv:2408.12755}
\author{FA Banach run packet}
\date{2026-06-23}

\begin{document}
\maketitle

\section*{Status}

This packet records a literature-implied full answer to Question
\texttt{quest:kpt} from Cuth--de Rancourt--Doucha,
\emph{Guarded Fraisse Banach spaces}, arXiv:2408.12755. The source asks:

\begin{quote}
Fix \(1\leq p,q<\infty\). Is the group \(\Iso(L_p(L_q))\), endowed with
the strong operator topology, extremely amenable? If not, does it have a
metrizable universal minimal flow?
\end{quote}

The answer is affirmative for every \(p,q\). Consequently the universal
minimal flow is a singleton.

The relation is not stated in the source and was not found as an explicit
later answer in the cheap index or exact web searches. It follows by combining
standard scalar \(L_p\)-isometry extreme amenability, Pestov--Schneider's
measurable-map theorem, and the Guerre--Raynaud structure theorem for
isometries of \(L_p(L_q)\).

\section*{Theorem}

Let \(1\leq p,q<\infty\), and let \(L_p(L_q)\) mean
\[
  L_p\bigl([0,1]; L_q([0,1])\bigr).
\]
Then \(\Iso(L_p(L_q))\), with the strong operator topology, is extremely
amenable. Equivalently, every continuous action of this group on a compact
Hausdorff space has a fixed point.

\section*{Proof}

We use three external facts.

First, the scalar isometry group \(\Iso(L_r)\), with the strong operator
topology, is extremely amenable for every \(1\leq r<\infty\): for \(r=2\) this
is the Gromov--Milman theorem for the infinite-dimensional orthogonal/unitary
group, and for \(r\neq 2\) this is the Giordano--Pestov theorem, as quoted in
the source discussion of Question \texttt{quest:kpt}.

Second, Pestov--Schneider prove that for every topological group \(G\),
\[
  \begin{aligned}
  G\text{ amenable}
  &\Longleftrightarrow
  L^0(G)\text{ amenable}\\
  &\Longleftrightarrow
  L^0(G)\text{ extremely amenable}\\
  &\Longleftrightarrow
  L^0(G)\text{ whirly amenable},
  \end{aligned}
\]
where \(L^0(G)\) is the group of strongly measurable maps
\([0,1]\to G\), modulo equality a.e., with convergence in measure.
Since extreme amenability implies amenability, this applies to
\(G=\Iso(L_q)\), giving that \(L^0(\Iso(L_q))\) is extremely amenable.

Third, for \(p\neq q\), the Guerre--Raynaud isometry theorem for vector-valued
\(L_p\)-spaces identifies the strong-operator-topology group
\(\Iso(L_p(L_q))\) with the natural measurable-field extension
\[
  1 \longrightarrow L^0(\Iso(L_q))
  \longrightarrow \Iso(L_p(L_q))
  \longrightarrow \Iso(L_p)
  \longrightarrow 1.
\]
Concretely, every surjective linear isometry of \(L_p([0,1];L_q)\) is a
base \(L_p\)-isometry together with a strongly measurable field of fiber
isometries of \(L_q\). The source paper itself invokes Guerre--Raynaud,
Theorem 5.1 and Proposition 1.6, in this form in the proof of Proposition
\texttt{prop:LpLqNotAuH}, where an isometry \(T\in \Iso(L_p(L_q))\) is
represented by an isometry \(S\in\Iso(L_p)\) and a measurable field
\((T_t)_{t\in[0,1]}\), \(T_t\in\Iso(L_q)\).

It remains only to use the elementary extension permanence of extreme
amenability. If \(N\triangleleft G\) is a closed normal subgroup, \(N\) is
extremely amenable, and \(G/N\) is extremely amenable, then \(G\) is extremely
amenable: for any compact \(G\)-flow \(K\), the fixed-point set \(K^N\) is
nonempty and compact; it is \(G\)-invariant, and the induced \(G/N\)-action on
\(K^N\) has a fixed point.

For \(p\neq q\), apply this lemma to the exact sequence above. The kernel
\(L^0(\Iso(L_q))\) is extremely amenable by Pestov--Schneider, and the quotient
\(\Iso(L_p)\) is extremely amenable by Gromov--Milman/Giordano--Pestov.
Therefore \(\Iso(L_p(L_q))\) is extremely amenable.

For \(p=q\), the Fubini isometry identifies
\[
  L_p([0,1];L_p([0,1])) \cong L_p([0,1]^2),
\]
which is isometric to the usual atomless separable \(L_p\)-space. Hence
\(\Iso(L_p(L_p))\cong\Iso(L_p)\) is extremely amenable by the same scalar
theorem.

This covers all \(1\leq p,q<\infty\).

\section*{Scope and Limitations}

The answer is full for Question \texttt{quest:kpt}: the ``if not'' alternative
does not arise, and the universal minimal flow is the one-point flow.

This packet does not address the other open questions in Section 7 of
arXiv:2408.12755. It also does not claim novelty of the component theorems;
the contribution is the direct identification that these known results answer
the source question.

\section*{Search and Provenance}

Cheap-index searches for arXiv:2408.12755, \emph{guarded Fraisse},
\(\Iso(L_p(L_q))\), and \emph{extremely amenable} found only the existing failed
attempt packet for the \(L_p(E)\) route and no promoted solution packet.
Exact web searches for the source title, guarded Fraisse Banach spaces,
\(L_p(E)\) guarded Fraisse, and \(\Iso(L_p(L_q))\) extremely amenable did not
surface a later explicit answer. Thus this is recorded as an
agent-identified literature implication rather than a later-paper
``already answered'' match.

\section*{References}

\begin{thebibliography}{9}

\bibitem{CDR}
M. Cuth, N. de Rancourt, and M. Doucha,
\emph{Guarded Fraisse Banach spaces}, arXiv:2408.12755.

\bibitem{GR}
S. Guerre and Y. Raynaud,
\emph{Sur les isometries de \(L^p(X)\) et le theoreme ergodique vectoriel},
Canad. J. Math. 40 (1988), no. 2, 360--391.

\bibitem{GM}
M. Gromov and V. D. Milman,
\emph{A topological application of the isoperimetric inequality},
Amer. J. Math. 105 (1983), no. 4, 843--854.

\bibitem{GP}
T. Giordano and V. Pestov,
\emph{Some extremely amenable groups related to operator algebras and
ergodic theory}, J. Inst. Math. Jussieu 6 (2007), no. 2, 279--315.

\bibitem{PS}
V. G. Pestov and F. M. Schneider,
\emph{On amenability and groups of measurable maps}, arXiv:1701.00281.

\end{thebibliography}

\end{document}
